% Target journal: Physical Review Research. The journal option prresearch exists only in
% REVTeX 4.2f and later; this manuscript is built with 4.2e, which has no such option, and
% APS applies the journal style at submission and production. The two-column layout of
% Physical Review Research is therefore set here by the class default plus the options below.
\documentclass[aps,10pt,superscriptaddress,floatfix,amsmath,amssymb,nofootinbib]{revtex4-2}

\usepackage[utf8]{inputenc}
\usepackage[T1]{fontenc}
\usepackage{amsmath,amssymb,amsthm}
\usepackage{url}
\usepackage[hidelinks,breaklinks]{hyperref}
\usepackage{bm}

\newcommand{\kb}{k_{\mathrm B}}
\newcommand{\Tlin}{T_{\mathrm{lin}}}
\newcommand{\Cinf}{C_\infty}
\newcommand{\Cthree}{C_{3}}
\newcommand{\Cstar}{C_{*}}
\newcommand{\Ccent}{C_{\mathrm{cent}}}
\newcommand{\Ccaus}{C_{\mathrm{caus}}}
\newcommand{\Ctwoone}{C_{21}}
\newcommand{\Cone}{C_{1}}
\newcommand{\Upre}{U}
\newcommand{\dstar}{d_{*}}
\newcommand{\drep}{d_{\mathrm{rep}}}
\newcommand{\mm}{\mathfrak m}
\newcommand{\Mmat}{\mathcal M}
\newcommand{\sS}{\mathcal S}
\newcommand{\Cgeom}{C_{\mathrm{geom}}}
\newcommand{\CLC}{C_{\mathrm{LC}}}
\newcommand{\Tceil}{T_{\mathrm{ceil}}}
\newcommand{\muS}{\mu^{*}}
\newcommand{\PROVED}{\textsc{proved}}
\newcommand{\CERT}{\textsc{certified}}
\newcommand{\NUM}{\textsc{numerical}}
\newcommand{\OPEN}{\textsc{open}}
\newcommand{\KNOWN}{\textsc{known}}
\newcommand{\DOI}[1]{\href{https://doi.org/#1}{\nolinkurl{#1}}}
\newcommand*{\PAPERDOI}{10.5281/zenodo.24411775}
\newcommand*{\SUPPLEMENTDOI}{10.5281/zenodo.24411779}

\newtheorem{theorem}{Theorem}
\newtheorem{lemma}{Lemma}
\newtheorem{corollary}{Corollary}
\newtheorem{proposition}{Proposition}
\newtheorem{definition}{Definition}
\newtheorem{remark}{Remark}

\begin{document}

\title{Certified ceilings for the Migdal--Eliashberg linear pairing onset:\\
one inequality chain from the glue moment to observable ratios}

\author{Davis Rens\homepage{ORCID 0000-0001-5839-5197}}
\affiliation{SuperMatics Inc.\\\href{https://orcid.org/0000-0001-5839-5197}{ORCID 0000-0001-5839-5197}}

\date{20 November 2026}

\begin{abstract}
We consolidate, in one notation and with the reduction proved in full, a chain of
rigorous upper bounds on the linear pairing onset temperature $\Tlin$ of a reciprocal
multichannel Migdal--Eliashberg model with channel-dependent bounded densities of states
and a common retarded glue. The bounds read $\kb\Tlin<C\sqrt{\rho}$, where $\rho$ is the
norm of the glue moment matrix, and they come in two strengths: a supremum bound
$C$ over the whole model class, and a strict per-model bound at a larger constant
$C'$ with a rational slack. For the centered class of normal branches the sharpest
supremum constant is $\Cthree=\sqrt{6811/5000}/(2\pi)=0.1857550244$ and the sharpest
strict constant is $\Cstar=\sqrt{341/250}/(2\pi)=0.1858777115$ with slack $9/6820$. For
the wider causal class defined by Harnack inequalities alone the sharpest strict constant
is $\Ccaus=\sqrt{45/32}/(2\pi)=0.1887345454$. The scalar Einstein limit gives the exact
constant $\Cinf=0.18272624777232952$, whose certified enclosure has relative width below
$6\times10^{-17}$; in a soft-glue limit with equitably coupled channels the coupled ceiling
is attained at $\Cinf$, and no method of the frozen-sequence type can produce a constant
below $0.185411$, so the coupled constants lie in an interval of certified width above the
exact value. A pairing-only Coulomb repulsion, at any pseudopotential, leaves the ceiling
intact while it destroys the optimality of the Einstein spectrum, and the two statements
are proved separately. The reduction from one edge inequality to the ceiling is proved
here in full. The finite part of every ceiling is an exact rational inequality verified
twice by independent programs; its statement and its record are reproduced. We then place
the ceilings against the literature: the Allen--Dynes strong-coupling coefficient $0.15$ is
$17.9\%$ below $\Cinf$, the Moussa--Cohen fixed-moment constant $0.183$ is $0.15\%$ above it, and
the coupling-resolved ceiling $C(\lambda)$ rises from $0.11462$ at $\lambda=1$ to $\Cinf$.
Measured transition temperatures of Pb, Nb, $\mathrm{MgB_2}$ and $\mathrm{LaH_{10}}$ sit at
$0.43$, $0.23$, $0.34$ and $0.60$ of the scalar ceiling, a ratio that is numerical
throughout. Retaining the first crossed skeleton, the answer is a dichotomy: on nested flat
sheets the crossed correction does not vanish in the adiabatic limit, on curved Fermi
surfaces it vanishes linearly in the adiabatic ratio up to a logarithm cubed, a fourth
power in two dimensions, and the
all-coupling all-momentum statement remains open. A limits section collects every object
that these bounds do not describe. No statement here identifies $\Tlin$ with the transition
temperature of a material, and no material is claimed to be a superconductor by these
bounds.
This preprint has not been peer reviewed.

\end{abstract}

\maketitle

\section{Introduction and summary}

The Migdal--Eliashberg equations are the working theory of phonon-mediated
superconductivity, and the empirical rule that a transition temperature follows the
electron-phonon coupling has been quantified by fits for sixty years
\cite{mcmillan,allendynes,moussa}. Fits are not bounds. The question addressed here is
the one a bound answers: how large can the pairing onset of the linearized theory be, as
a function of quantities that a material calculation provides, uniformly over a class of
spectra and of densities of states?

Two numbers measure a coupling. The first is $\lambda$, the mass of the glue. The second
is $M=\lambda\langle\omega^{2}\rangle$, the second moment of the Eliashberg function
weighted by the coupling, equivalently $\lambda$ times the mean square phonon energy. The
bounds in this paper are formulated with the matrix of glue moments of a multichannel
model and with its operator norm $\rho$; in the single-channel case $\rho=M$, so every
constant below can be read as a coefficient multiplying $\sqrt{M}$, and then converted to
kelvin by the mean phonon energy.

The physical content is a ceiling of order $0.18\sqrt{\lambda}\,\omega$ with a certified
coefficient. Four consequences follow, and each is stated with its own status label in
what follows. First, the coefficient is not a fit: it is the exact value of a
strong-coupling limit, enclosed to seventeen decimal places, and the Allen--Dynes
strong-coupling coefficient $0.15$ lies $17.9\%$ below it, so the fit is not merely unproved but
strictly conservative in that regime. Second, for a single channel the ceiling saturates
as $\lambda\to\infty$, which is the statement that strong coupling alone cannot push the
linear onset above a fixed multiple of the phonon energy scale. Third, in the multichannel
problem the ceiling depends on the glue matrix through its norm and not through its
eigenvalues individually, so a bipartite glue whose moment matrix has a negative
eigenvalue is covered with the same constant. Fourth, the coupled constant is pinned from
below as well: in a soft-glue limit with equitably coupled channels the sharp scalar value
$\Cinf$ is attained, so the width of the interval that contains the coupled supremum is
the width between $\Cinf$ and the certified constant, and not a matter of opinion about
which channel structure is optimal.

\subsection{What this paper contains}

This paper is a consolidated record. Its statements are those of eleven released and
independently reviewed records
\cite{finiteband,gamma,univglue,optimum,finiteenergy,anisosharp,msvertex,cutoff,slowphonon,limits,capstone},
seven of which form the ceiling chain of Secs.~\ref{sec:model} to~\ref{sec:ceilings}. The
paper makes four contributions of its own, all of them organizational or expository:

\begin{enumerate}
\item[\rm(i)] It states the model, the two classes of normal branches and the onset
temperature once, in one notation, in Sec.~\ref{sec:model}.
\item[\rm(ii)] It proves the reduction from a single edge inequality to the ceiling in full
(Lemma~\ref{lem:reserve} and Corollary~\ref{cor:strict}). The published records prove it
too; reproducing it here makes the dependence of the final constant on exactly one
finite statement explicit.
\item[\rm(iii)] It places every constant in one ordered table with the reserve that
produces it, the exact rational slack it carries, and the certification record behind it
(Tables~\ref{tab:ceil} and~\ref{tab:reserve}), and it sets the constants against the
literature and against four sets of measured transition temperatures
(Secs.~\ref{sec:sharp} to~\ref{sec:lit}).
\item[\rm(iv)] It records the one place where the coupled problem is closed: the soft-glue
limit with equitably coupled channels, where the ceiling is attained at $\Cinf$
(Theorem~\ref{thm:softlimit}). This records a released statement in the notation of the
general model, and it is the reason the abstract interval $[\Cinf,\Cthree]$ is a statement
about a limit rather than a guess.
\end{enumerate}

No new analytic inequality is claimed, and the exact-arithmetic brackets of
Tables~\ref{tab:clambda} and~\ref{tab:mat} are presented here for the first time, with their
programs and digitization records in the supplement. Where a statement belongs to an
earlier record, the record
is identified by its persistent identifier, and the label attached to the statement is the
label the record carries after independent review. The labels are: \PROVED{} means a
human-reviewed analytic argument combined with exact rational arithmetic that an
independent program reproduces; \CERT{} means the finite exact statement on its own, apart
from the analysis that applies it; \NUM{} means floating-point evidence without an
enclosure; \OPEN{} means undecided. Nothing in this paper carries the label \PROVED{} by
itself except the reduction of Sec.~\ref{sec:reserve}, which is elementary.

\subsection{The result in one line}

For every model of the centered class of Sec.~\ref{sec:model},
\begin{equation}\label{eq:headline}
\kb\Tlin\;<\;\Cstar\sqrt{\rho},\qquad
\Cstar=\frac{\sqrt{341/250}}{2\pi}=0.1858777114719168\ldots,
\end{equation}
and any constant $C>\Cthree$ with $\Cthree=\sqrt{6811/5000}/(2\pi)=0.1857550244$ is a
valid ceiling for every model of the class, whereas $\Cthree$ itself is guaranteed only in
the non-strict form $\kb\Tlin\le\Cthree\sqrt{\rho}$. That is the difference between the
strict and the supremum column of Table~\ref{tab:ceil}. The exact scalar constant is $\Cinf=0.18272624777232952$, so the
coupled constant lies in the interval $[\Cinf,\Cthree]$, of relative width $1.66\%$; the
lower endpoint is attained in the soft-glue limit of Theorem~\ref{thm:softlimit}, and
whether a single model at finite moment attains it is open. For the causal class the
corresponding numbers are $\Ccaus=\sqrt{45/32}/(2\pi)=0.1887345454$ strict, and the
supremum bound $\sqrt{281/200}/(2\pi)=0.1886506447$, which is $1.6\%$ above $\Cthree$.

\section{Model, class of branches, and onset}\label{sec:model}

The model is the reciprocal multichannel Eliashberg model of
Refs.~\cite{univglue,optimum,finiteenergy,anisosharp}. We restate it because every
constant below is a statement about this object and about nothing else.

\subsection{Channels, densities of states and glue}

The channel space is a finite set $X$ with positive height weights $\eta_x$. Channel $x$
carries a measurable density of states $0\le N_x(\xi)\le\eta_x$ whose peak height is the
weight of the channel, $N_{\max,x}=\eta_x$, and whose normalized profile
$\nu_x=N_x/N_{\max,x}$ is even and valued in $[0,1]$. The vertices are energy independent
and radial, so a channel is a band and not a momentum direction.

The glue is a family of reciprocal nonnegative measures
$d\mm_{xy}(E)=d\mm_{yx}(E)\ge0$ on energies $E\in(0,\infty)$, with moment matrix
$\Mmat_{xy}=\int d\mm_{xy}(E)$. The matrix acts on $\ell^{2}(\eta)$ by
\begin{equation}\label{eq:moment}
(\Mmat h)_x=\sum_y\eta_y\Mmat_{xy}h_y,\qquad
\rho=\|\Mmat\|_{\ell^{2}(\eta)}\in(0,\infty).
\end{equation}
Entrywise nonnegativity is required and positive semidefiniteness of $\Mmat$ is not: a
bipartite glue, one that couples two groups of channels without coupling within a group,
has a moment matrix with a negative eigenvalue, so the relaxation is physical. Each measure
is supported on positive energies and has a finite moment, and at each tested temperature
its zero-transfer mass is finite. In a single channel with an Einstein line of energy
$\Omega$ and coupling $\lambda$ one has $\Mmat=\lambda\Omega^{2}$ and $\rho=\lambda
\Omega^{2}=M$. For a general spectrum, $\rho\le\lambda\max_x\langle\omega^{2}\rangle_x$,
with equality when all channels share one spectrum.

\subsection{A positive reflected normal branch}

Set $\tau=\kb T$ and $\omega_n=\pi\tau w_n$ with $w_n=2n+1$. Pairing is tested on a
positive reflected normal branch, that is, on an actual solution of the matched
imaginary-axis normal equation. Its normal width at the $n$th reflected frequency, in
units in which the free width is $w_n$, is
\begin{equation}\label{eq:width}
s_{n,x}=w_n+\frac{1}{(2\pi\tau)^{2}}\sum_y\eta_y\int\sum_{m\ge0}p_{nm}(t)\,b_{m,y}\,
d\mm_{xy}(E)\ \ge\ w_n,
\end{equation}
and the pairing weights are the Poisson averages of the profile,
\begin{equation}\label{eq:weights}
b_{n,x}=\int_{-\infty}^{\infty}\nu_x(\xi)\,
\frac{\pi\tau s_{n,x}}{(\pi\tau s_{n,x})^{2}+\xi^{2}}\,\frac{d\xi}{\pi}\in[0,1].
\end{equation}
An extra allowed normal mass only increases the width in \eqref{eq:width}, so nothing
below depends on assuming the minimal normal kernel. The kernel factors are
\begin{equation}\label{eq:kernel}
k_j(t)=\frac{1}{t+j^{2}},\quad
p_{nm}=k_{|n-m|}-k_{n+m+1}\ge0,\quad
S_{nm}=k_{|n-m|}+k_{n+m+1},\qquad t=\Bigl(\frac{E}{2\pi\tau}\Bigr)^{2},
\end{equation}
the factor $(2\pi\tau)^{2}$ having been removed from the physical transfer.

\subsection{Two classes of branches}

The widths \eqref{eq:width} are not free. In the \emph{centered class} the ratio
$s_{n,x}/w_n$ is nonincreasing and $w_ns_{n,x}$ is nondecreasing in $n$, which gives the
width cone
\begin{equation}\label{eq:cone}
w_m^{-1}\ \le\ \frac{s_{m,x}}{s_{0,x}}\ \le\ w_m,\qquad m\ge1.
\end{equation}
The two centered Poisson kernels at heights whose ratio lies in $[w_m^{-1},w_m]$ have a
pointwise ratio in the same interval, so integrating both the profile and its complement
in \eqref{eq:weights} gives the four Harnack inequalities
\begin{equation}\label{eq:harnack}
0\le b_{m,x}\le1,\qquad
\frac{b_{0,x}}{w_m}\le b_{m,x}\le w_mb_{0,x},\qquad
w_mb_{0,x}-(w_m-1)\le b_{m,x}\le\frac{b_{0,x}+w_m-1}{w_m}.
\end{equation}
The \emph{causal class} is the wider class in which \eqref{eq:harnack} is retained and
\eqref{eq:cone} is dropped, the weights being those of a channel-wise causal radial
propagator. Every centered branch is causal, and the constants for the causal class are
correspondingly larger. The cone \eqref{eq:cone} was established for the centered class in
Refs.~\cite{univglue,optimum}, and the Harnack form \eqref{eq:harnack} for the causal class
in Ref.~\cite{finiteenergy}.

\begin{remark}[What is assumed and what is not]\label{rem:assume}
At each tested temperature the frequency rows of the normal kernel, including the
zero-transfer term, are assumed to have a finite bound $R(\tau)$; no bound uniform in
$\tau$ is needed anywhere. The normal kernel is allowed to exceed the pairing transfer by
an odd-folded nonnegative excess, and a positive semidefinite static pairing-only
subtraction is harmless on the matched branch. These relaxations are those of
Refs.~\cite{optimum,anisosharp} and are used in the reduction of
Sec.~\ref{sec:reserve} only through the existence of $R(\tau)$.
\end{remark}

\subsection{The stability form and the onset temperature}

For a finite-support real profile $v_{n,x}$ put $g_{n,x}=\sqrt{w_nb_{n,x}}\,v_{n,x}$.
After folding the frequency axis the normal-state stability form is
\begin{equation}\label{eq:form}
Q(v)=\sum_x\eta_x\|g_x\|^{2}
+\frac{1}{2(2\pi\tau)^{2}}\sum_{x,y}\eta_x\eta_y\int I_{xy}(t)\,d\mm_{xy}(E),
\end{equation}
where, with $b_n=b_{n,x}$, $c_m=b_{m,y}$, $X_n=v_{n,x}$ and $Y_m=v_{m,y}$,
\begin{equation}\label{eq:Ixy}
I_{xy}(t)=\sum_{n,m\ge0}b_nc_m\Bigl[p_{nm}(X_n^{2}+Y_m^{2})-2S_{nm}X_nY_m\Bigr].
\end{equation}

\begin{definition}[Linear onset]\label{def:onset}
$\Tlin$ is the supremum of the temperatures at which $Q$ fails to be positive definite on
finitely supported real profiles of a single symmetry sector; for a model with several
sectors it is the largest of the sector onsets. It is an energy-valued linear onset
temperature. It is not the transition temperature of an exact Hamiltonian, and it is not
measured.
\end{definition}

Real and imaginary parts of a trial function decouple in \eqref{eq:form}, because
$I_{xy}$ is a real quadratic form in $(X,Y)$, so real trials suffice throughout.

\section{The reserve lemma: from one edge inequality to the ceiling}\label{sec:reserve}

Everything analytic in this paper is one lemma. It says that a bound on the stability form
for a \emph{single pair of channels at a single energy}, valid uniformly over the admissible
weights, the imbalance between the two channels and the energy, implies a bound on $Q$ for
any number of channels and any glue, with a constant that depends on the pair bound only.
The proof is the positive-resolvent argument of Refs.~\cite{finiteenergy,anisosharp},
reproduced in full.

For channels $x,y$, real weights $(b_n)_{n\ge0}$ and $(c_m)_{m\ge0}$ and $r>0$ put
\begin{equation}\label{eq:Br}
B_r=\sum_{n\ge0}w_n\,r^{-1}b_nX_n^{2}+\sum_{m\ge0}w_m\,r\,c_mY_m^{2}
=r^{-1}\|g_x\|^{2}+r\|g_y\|^{2},
\end{equation}
where $\|g_x\|^{2}=\sum_nw_nb_nX_n^{2}$ and $\|g_y\|^{2}=\sum_mw_mc_mY_m^{2}$.

\begin{proposition}[edge inequality]\label{prop:edge}
Fix a class of branches and a real number $0<d<1$. Suppose that for every $t>0$, every
$r>0$, every pair of weight sequences $(b_n),(c_m)$ admissible in that class at any
temperature, and every real finitely supported $X,Y$,
\begin{equation}\label{eq:edge}
B_r+\tfrac12 I_{xy}(t)\ \ge\ d\,B_r .
\end{equation}
Then $d$ is called an \emph{edge reserve} for the class. The reserve carries no dependence
on $t$, on $r$ or on the weights.
\end{proposition}

\begin{lemma}[reserve lemma]\label{lem:reserve}
Let $d$ be an edge reserve for a class. Then for every model of that class and every
$\tau>0$,
\begin{equation}\label{eq:Qbound}
Q(v)\ \ge\ \Bigl[1-\frac{2(1-d)\rho}{(2\pi\tau)^{2}}\Bigr]\sum_x\eta_x\|g_x\|^{2}
\end{equation}
for every finitely supported real profile, and consequently
$\kb\Tlin\le\dfrac{\sqrt{2(1-d)\rho}}{2\pi}$.
\end{lemma}

\begin{proof}
Choose $R'>\rho$ and put
$h=(I-\Mmat/R')^{-1}\mathbf 1=\sum_{j\ge0}(\Mmat/R')^{j}\mathbf 1$. Every entry of $\Mmat$
is nonnegative and $\mathbf 1>0$, so $h>0$ entrywise and
\begin{equation}\label{eq:hbound}
\Mmat h=R'(h-\mathbf 1)\ \le\ R'h\quad\text{entrywise,}
\end{equation}
and no positivity of the quadratic form of $\Mmat$ is used anywhere.

Apply \eqref{eq:edge} with $r=r_{xy}=h_x/h_y$ for every ordered pair $(x,y)$ and integrate
against $d\mm_{xy}$. By \eqref{eq:Br} the first term is
$(h_y/h_x)\|g_x\|^{2}+(h_x/h_y)\|g_y\|^{2}$, so
\begin{equation}\label{eq:step1}
Q(v)\ \ge\ \sum_x\eta_x\|g_x\|^{2}
-\frac{1-d}{(2\pi\tau)^{2}}\sum_{x,y}\eta_x\eta_y\Mmat_{xy}
\Bigl(\frac{h_y}{h_x}\|g_x\|^{2}+\frac{h_x}{h_y}\|g_y\|^{2}\Bigr).
\end{equation}
By reciprocity $\Mmat_{xy}=\Mmat_{yx}$ the two terms in the bracket give the same double
sum. That sum equals
\begin{equation}\label{eq:step2}
\sum_x\frac{\eta_x\|g_x\|^{2}}{h_x}\sum_y\eta_y\Mmat_{xy}h_y
\;=\;\sum_x\frac{\eta_x\|g_x\|^{2}}{h_x}(\Mmat h)_x
\;\le\;R'\sum_x\eta_x\|g_x\|^{2}
\end{equation}
by \eqref{eq:hbound}, and it is nonnegative. Hence
\begin{equation}\label{eq:gap}
Q(v)\ \ge\ \Bigl[1-\frac{2(1-d)R'}{(2\pi\tau)^{2}}\Bigr]\sum_x\eta_x\|g_x\|^{2}.
\end{equation}
The left side does not depend on $R'$, so we may let $R'\downarrow\rho$.

Now fix $\tau$ with $(2\pi\tau)^{2}>2(1-d)\rho$ and put
$\beta(\tau)=1-2(1-d)\rho/(2\pi\tau)^{2}>0$. Letting $R'\downarrow\rho$ in \eqref{eq:gap}
gives $Q(v)\ge\beta(\tau)\sum_x\eta_x\|g_x\|^{2}$ for every $v$. The normal widths are
bounded in terms of the finite row bound $R(\tau)$ of Remark~\ref{rem:assume} by
$s_{n,x}\le w_n+R(\tau)\le(1+R(\tau))w_n$ at this temperature, so the Hilbert norm of the
pairing problem satisfies
\begin{equation}\label{eq:norm}
\sum_x\eta_x\bigl\|(s_xb_x)^{1/2}v_x\bigr\|^{2}
=\sum_x\eta_x\sum_n s_{n,x}b_{n,x}v_{n,x}^{2}
\le(1+R(\tau))\sum_x\eta_x\|g_x\|^{2}.
\end{equation}
Therefore
\begin{equation}\label{eq:final}
Q(v)\ \ge\ \frac{\beta(\tau)}{1+R(\tau)}\sum_x\eta_x\bigl\|(s_xb_x)^{1/2}v_x\bigr\|^{2}
\end{equation}
for every $v$, and $Q$ is strictly positive definite at $\tau$. No such $\tau$ is a
failure temperature, so $\kb\Tlin\le\sqrt{2(1-d)\rho}/(2\pi)$. Disconnected instability
sets are included, since the argument is applied at each $\tau$ and never along a branch.
\end{proof}

\begin{corollary}[strict bound from two reserves]\label{cor:strict}
Let $d$ and $d'$ be edge reserves for the same class with $d>d'$. Then for every model of
the class,
\begin{equation}\label{eq:strict}
\kb\Tlin\le\frac{\sqrt{2(1-d)\rho}}{2\pi}<\frac{\sqrt{2(1-d')\rho}}{2\pi},
\end{equation}
and at the larger temperature $\tau'=\sqrt{2(1-d')\rho}/(2\pi)$ the form itself satisfies
\begin{equation}\label{eq:slack}
Q(v)\ \ge\ \frac{d-d'}{1-d'}\sum_x\eta_x\|g_x\|^{2}
\end{equation}
before the width conversion \eqref{eq:norm}.
\end{corollary}

\begin{proof}
The first statement is Lemma~\ref{lem:reserve} applied twice, using $d>d'$. For the second,
insert $\rho=(2\pi\tau')^{2}/(2(1-d'))$ into \eqref{eq:gap} with $R'\downarrow\rho$:
\[
1-\frac{2(1-d)\rho}{(2\pi\tau')^{2}}
=1-\frac{1-d}{1-d'}=\frac{d-d'}{1-d'}>0 .
\]
\end{proof}

\begin{remark}[Why two reserves are needed]\label{rem:tworeason}
A single reserve $d$ gives a ceiling, but not a strict inequality at a particular model:
stability at one temperature is not by itself stability in a neighbourhood of it, and no
monotonicity of the onset in temperature is used anywhere in this paper. The strict
statement is therefore read off a second reserve $d'<d$ of the same inequality, whose
survival at a temperature strictly above the ceiling is what closes the argument, and the
rational margin \eqref{eq:slack} records the slack. The two reserves are different
thresholds of the same edge inequality, not a chain of increasingly strong hypotheses: a
smaller reserve does not by itself yield a smaller ceiling, and each strict constant is
computed from its own covering. This is the logical structure of
Refs.~\cite{finiteenergy,anisosharp}, and the slacks $1/2050$, $1/1125$, $9/6820$ and
$1/754$ in Table~\ref{tab:reserve} are the numbers that appear in the released statements.
\end{remark}

\begin{remark}[The role of connectedness]\label{rem:connect}
The proof never uses a branch. It uses one reflected normal branch at one temperature,
finite rows $R(\tau)$, and the edge inequality. A model whose instability set is
disconnected, or which has no continuous branch at all, is covered as long as a positive
reflected normal branch exists at the tested temperatures.
\end{remark}

\section{The certified finite part: edge inequalities and records}\label{sec:finite}

By Lemma~\ref{lem:reserve} the whole analytic problem is reduced to the single inequality
\eqref{eq:edge}. That inequality is finite dimensional but not small: it involves all
frequencies, all energies and all admissible weights. The released records reduce it
further, by an enclosure of the head weight against a finite set of tail frequencies in a
fixed number of stages, to a finite set of rational inequalities on rectangles, and then
verify those in exact arithmetic. We state the outcome and the record, and we describe the
mechanism, because the mechanism is what decides which reserve belongs to which class.

\subsection{The reserves}

\begin{theorem}[edge reserves]\label{thm:reserves}
The following numbers are edge reserves in the sense of Proposition~\ref{prop:edge}. In
each pair the first member is the reserve that is verified with equality allowed, so it
yields a supremum statement and the non-strict ceiling, and the second is a smaller reserve,
verified with the slack left in, which \eqref{eq:slack} converts into a strict ceiling at a
somewhat larger constant:
\begin{enumerate}
\item[\rm(a)] $d=317/1000$ and $d'=19/60$, centered class, Ref.~\cite{finiteenergy};
\item[\rm(b)] $d=119/400$ and $d'=19/64$, causal class, Ref.~\cite{finiteenergy};
\item[\rm(c)] $d=247/1000$ and $d'=123/500$, causal class, Ref.~\cite{optimum};
\item[\rm(d)] $d=3189/10000$ and $d'=159/500$, centered class, Ref.~\cite{anisosharp}.
\end{enumerate}
Status: \PROVED{} for each pair, conditionally on the exact finite statement of
Sec.~\ref{sec:finite}, which is \CERT{}: it was verified twice by independently written
exact-arithmetic programs and, in its finite core, in Lean~4.35.0. The analytic part
(the Harnack and Poisson rearrangement, the width cone, the cascade, the head criterion and
the exterior regions) was reviewed in Refs.~\cite{finiteenergy,anisosharp} and is not
reproduced here; the reduction that uses it is Lemma~\ref{lem:reserve}, which is complete.
\end{theorem}

\subsection{The finite statement and what it certifies}

\begin{proposition}[finite core, statement]\label{prop:finitecore}
Fix a class and a reserve $d$. Let $(r,t)$ run over the compact region
$[1,3]\times[0,1/2]$, the primary enclosure of Refs.~\cite{finiteenergy,anisosharp}, or over
the symmetry-free region $[1/3,3]\times[0,1/2]$ used by the second program of
Ref.~\cite{anisosharp}. Split the head weight interval into sub-boxes, and in
Ref.~\cite{finiteenergy} the head weights $b_0,c_0$ into nine sub-intervals each, and
complete the head frequency against twenty tail frequencies in three stages: Stage~0 pairs
each tail mode with the head, Stage~1 pairs the first tail mode with the higher tail modes,
and Stage~2 pairs the second tail mode with the higher ones. The finite statement is that on
every leaf of this decomposition the exact rational expression obtained by dropping no
positive term and bounding every negative one by the retained completions is nonnegative,
with equality allowed for the reserve $d$ and with the rational slack \eqref{eq:slack} kept
for the reserve $d'$. Outside the compact region a single cruder argument is used, with the
share $\alpha_{\mathrm{ext}}=617/1000$ left by the full Hankel bound in
Ref.~\cite{finiteenergy}; in Ref.~\cite{anisosharp} the exterior region uses neither the
cascade nor the head minimisation.
Status: \CERT{}.
\end{proposition}

The primary covering of Ref.~\cite{anisosharp} needs $367$ leaves at depth $24$ and the
independent one $1036$ leaves at depth $13$ for the centered reserve $3189/10000$, against
$159$ and $466$ leaves for the reserve $159/500$ of its first version. The primary program
uses the symmetry $r\leftrightarrow r^{-1}$ to cover an area of one and halves the $t$
interval while the $r$ interval is still wide, so that the tail frequencies are resolved
before the head weight; the second program does not use the symmetry and covers the full
domain of area $4/3$ by four-way subdivision. The reserve that the search certified is the
largest it reached inside its depth budget, not the largest the scheme can reach: the
shallower caps of the first version certified $3188/10000$ and the caps used here certified
$3189/10000$, and the first attempt at $3189/10000$ failed only at those caps, so that
raising a cap is bookkeeping and weakens no inequality. The next grid point $3190/10000$
exhausted the remaining time budget and is left unresolved, neither rejected nor shown to be
impossible. Each guarded run corrupts a
production quantity and must exit with status $1$ and a line beginning \texttt{REJECT}; the
released records report fourteen expected guarded mutant failures for Ref.~\cite{optimum},
thirteen and eight mutated runs for the two programs of Ref.~\cite{finiteenergy}, and ten
guarded runs for each program of Ref.~\cite{anisosharp}. The tightest leaves are recorded
explicitly, for example
$(r,t)\in[135/128,271/256]\times[0,1/256]$ with head box $I_6\times I_5$ and certified slack
$0.000078371512$ for the centered reserve $317/1000$, and
$(r,t)\in[131/128,33/32]\times[0,1/256]$ with certified slack $0.000003769886$ for the
causal reserve $19/64$, both in Ref.~\cite{finiteenergy}.

\subsection{The mechanism, and why the reserves are ordered as they are}

Three ingredients decide the number.

The \emph{head} frequency carries the largest weight and is completed against the tail.
Completing the head against each tail with a static share of the tail mass stalls: at the
least favourable imbalance the first tail frequency falls short of its stability threshold
$2/9$ once the reserve approaches $0.30$, because the static share satisfies
$\alpha_{0,1}\le0.634$ and $r\alpha_{0,1}\le0.2113<2/9$ at $r=1/3$ (Ref.~\cite{finiteenergy},
discussion of the Stage-1 mechanism). That is why the argument without the first stage,
which is the argument of Ref.~\cite{optimum}, stops at the reserve $247/1000$ and cannot
approach $0.30$, whereas the staged argument of Ref.~\cite{finiteenergy} reaches $317/1000$.

The \emph{cascade} removes the stall by stabilising the first tail mode against the higher
tail modes before the head is touched, which raises its effective normal mass. One exact
difference-of-squares identity carries all stages, written
\begin{equation}\label{eq:cascade}
S_{mk}^{2}-p_{mk}^{2}=\frac{4p_{mk}}{w_mw_k}
\end{equation}
in the notation of Ref.~\cite{finiteenergy}, so a single monotonicity statement suffices,
and every gain is delivered to the head completion of its own frequency, with the head
weight as the only divisor and no gain reused. With the cascade the centered reserve rises
from $19/60$, the strict value of Ref.~\cite{finiteenergy}, through $159/500$, the value of
the first version of Ref.~\cite{anisosharp}, to $3189/10000$; the causal reserve, which
cannot use the cone, stays at $19/64$.

The \emph{class} decides whether the widths may be used. In the centered class the cone
\eqref{eq:cone} is available and gives the sharper lower envelope on the tail weights; in
the causal class only \eqref{eq:harnack} is retained, with reserves $119/400$ and $19/64$
in the third row of Table~\ref{tab:reserve}, which is why the two causal constants sit
above the corresponding centered constants of the same record: by $1.42\%$ in the
supremum form and by $1.44\%$ in the strict form.

\section{The ceilings}\label{sec:ceilings}

\begin{theorem}[ceilings]\label{thm:ceil}
For every model of the stated class, with $\rho$ the norm \eqref{eq:moment} of the glue
moment matrix:
\begin{enumerate}
\item[\rm(a)] \emph{Centered class, supremum.}
The supremum $\sup\kb\Tlin/\sqrt\rho$ is at most $\Cthree$, where
$\Cthree=\sqrt{6811/5000}/(2\pi)=0.1857550243988$, the supremum being over all models of the
class with finitely many channels, all finite positive glue energies and all admissible
centered branches.
\item[\rm(b)] \emph{Centered class, strict.}
$\kb\Tlin<\Cstar\sqrt\rho$ with $\Cstar=\sqrt{341/250}/(2\pi)=0.1858777114719$, the bare
slack at $\tau=\Cstar\sqrt\rho$ being the rational $9/6820$.
\item[\rm(c)] \emph{Causal class, strict.}
$\kb\Tlin<\Ccaus\sqrt\rho$ with $\Ccaus=\sqrt{45/32}/(2\pi)=0.1887345454$, the bare slack at
$\tau=\Ccaus\sqrt\rho$ being $1/1125$; and $\kb\Tlin<\Ctwoone\sqrt\rho$ with
$\Ctwoone=\sqrt{377/250}/(2\pi)=0.1954433070$, slack $1/754$.
\item[\rm(d)] \emph{Centered class, earlier constants.}
$\kb\Tlin<\Ccent\sqrt\rho$ with $\Ccent=\sqrt{41/30}/(2\pi)=0.1860593214$, slack $1/2050$;
$\kb\Tlin\le\sqrt{683/500}/(2\pi)\,\sqrt\rho=0.1860139355\sqrt\rho$, the corresponding
supremum bound of the same record; and the two earlier coefficients of
Refs.~\cite{univglue,optimum}, $\kb\Tlin\le\Cone\sqrt\rho$ with
$\Cone=\sqrt{413/250}/(2\pi)=0.2045620904$, obtained by an additive first-frequency
comparison without a recorded slack, and $\kb\Tlin\le\Upre\sqrt\rho$ with
$\Upre=\sqrt{6466219/3576050}/(2\pi)=0.2140146706$ for the centered class of the earlier
note. The last two are not consequences of an edge reserve, so their status is that of their
records and not that of Theorem~\ref{thm:reserves}.
\item[\rm(e)] \emph{Causal class, supremum.}
$\sup\kb\Tlin/\sqrt\rho\le\sqrt{281/200}/(2\pi)=0.1886506447$ for the causal class, from
the reserve $119/400$ of Ref.~\cite{finiteenergy}, and
$\le\sqrt{753/500}/(2\pi)=0.1953136596$ from the reserve $247/1000$ of
Ref.~\cite{optimum}.
\end{enumerate}
Status: \PROVED{} in the sense of Theorem~\ref{thm:reserves}: analytic reduction complete
here, finite part certified in exact arithmetic, statements released and independently
reviewed in Refs.~\cite{univglue,optimum,finiteenergy,anisosharp}.
\end{theorem}

\begin{proof}
Apply Lemma~\ref{lem:reserve} with the reserves of Theorem~\ref{thm:reserves} and
Corollary~\ref{cor:strict} where two reserves for the same class are available, using
$2(1-3189/10000)=6811/5000$, $2(1-159/500)=341/250$, $2(1-317/1000)=683/500$,
$2(1-19/60)=41/30$, $2(1-119/400)=281/200$, $2(1-19/64)=45/32$,
$2(1-247/1000)=753/500$ and $2(1-123/500)=377/250$, which is the content of \rm(a) to \rm(c)
and of \rm(e). The slacks of \rm(b), \rm(c) and \rm(d) are the quotients \eqref{eq:slack},
$(3189/10000-159/500)/(341/500)=9/6820$,
$(119/400-19/64)/(45/64)=1/1125$,
$(317/1000-19/60)/(41/60)=1/2050$ and $(247/1000-123/500)/(377/500)=1/754$.
Items \rm(d) for $\Cone$ and $\Upre$ are not of the reserve form and are imported from their
records; they appear here so that the table is ordered, not because this proof produces
them.
\end{proof}

\begin{remark}[Ordering]\label{rem:ordering}
Theorem~\ref{thm:ceil} gives a chain of strict inequalities between explicit numbers,
\begin{equation}\label{eq:chain}
\Cinf<\Cthree<\Cstar<\frac{\sqrt{683}}{2\pi\sqrt{500}}<\Ccent<\frac{\sqrt{281/200}}{2\pi}
<\Ccaus<\frac{\sqrt{753/500}}{2\pi}<\Ctwoone<\Cone<\Upre ,
\end{equation}
each link after the first being a comparison of rationals under a square root. The first
constant is the exact scalar value of Sec.~\ref{sec:sharp}, and the inequality
$\Cinf<\Cthree$ is read off the certified enclosure of that section. The two smallest
coupled constants are the sharpest known for the centered class, and they lie $1.66\%$ and
$1.72\%$ above the scalar value.
\end{remark}

\begin{table}[h]
\caption{The ceilings, ordered. The last column is the relative excess over the exact scalar
constant $\Cinf=0.18272624777232952$. Status \emph{supremum} means that the number bounds
$\sup\kb\Tlin/\sqrt\rho$ over the class, so for an individual model the statement is
non-strict; \emph{strict} means $\kb\Tlin<C\sqrt\rho$ for every model of the class, with the
rational slack of Table~\ref{tab:reserve}; \emph{earlier} means superseded. The records
behind the rows are, in order of appearance, Ref.~\cite{gamma}, Ref.~\cite{anisosharp} for
$\Cthree$ and $\Cstar$, Ref.~\cite{finiteenergy} for the centered pair and the causal pair,
Ref.~\cite{optimum} for $\Ctwoone$ and $\Cone$, and Ref.~\cite{univglue} for $\Upre$. Labels
follow Sec.~\ref{sec:model}: $\Cinf$ is \PROVED{}, the reserves are \PROVED{} subject to a
\CERT{} certificate, and the decimals follow by elementary arithmetic.\label{tab:ceil}}
\small
\begin{tabular}{@{}llllll@{}}
\hline
constant & class & status & exact form & decimal & excess over $\Cinf$ \\
\hline
$\Cinf$ & scalar & exact & transcendental & $0.1827262478$ & \\
$\Cthree$ & centered & supremum & $\sqrt{6811/5000}/(2\pi)$ & $0.1857550244$ & $1.658\%$ \\
$\Cstar$ & centered & strict & $\sqrt{341/250}/(2\pi)$ & $0.1858777115$ & $1.725\%$ \\
$\sqrt{683/500}/(2\pi)$ & centered & supremum & $\sqrt{683/500}/(2\pi)$ & $0.1860139355$ & $1.799\%$ \\
$\Ccent$ & centered & strict & $\sqrt{41/30}/(2\pi)$ & $0.1860593214$ & $1.824\%$ \\
$\sqrt{281/200}/(2\pi)$ & causal & supremum & $\sqrt{281/200}/(2\pi)$ & $0.1886506447$ & $3.242\%$ \\
$\Ccaus$ & causal & strict & $\sqrt{45/32}/(2\pi)$ & $0.1887345454$ & $3.288\%$ \\
$\sqrt{753/500}/(2\pi)$ & causal & supremum & $\sqrt{753/500}/(2\pi)$ & $0.1953136596$ & $6.889\%$ \\
$\Ctwoone$ & causal & strict & $\sqrt{377/250}/(2\pi)$ & $0.1954433070$ & $6.960\%$ \\
$\Cone$ & causal & earlier & $\sqrt{413/250}/(2\pi)$ & $0.2045620904$ & $11.95\%$ \\
$\Upre$ & centered & earlier & $\sqrt{6466219/3576050}/(2\pi)$ & $0.2140146706$ & $17.12\%$ \\
\hline
\end{tabular}
\end{table}

\begin{table}[h]
\caption{Reserves and slacks. A reserve $d$ gives a ceiling at $\sqrt{2(1-d)}/(2\pi)$; two
reserves $d>d'$ for the same class give the strict statement at $\sqrt{2(1-d')}/(2\pi)$ with
the bare slack $(d-d')/(1-d')$ of Corollary~\ref{cor:strict}. Leaves and guarded runs are
the certification records of the named preprints.\label{tab:reserve}}
\small
\begin{tabular}{@{}llllll@{}}
\hline
class & reserve $d$ & reserve $d'$ & slack & constant(s) & record \\
\hline
centered & $3189/10000$ & $159/500$ & $9/6820$ & $\Cthree$, $\Cstar$ & 24410276 \\
centered & $317/1000$ & $19/60$ & $1/2050$ & $\sqrt{683/500}/(2\pi)$, $\Ccent$ & 24380829 \\
causal & $119/400$ & $19/64$ & $1/1125$ & $\sqrt{281/200}/(2\pi)$, $\Ccaus$ & 24380829 \\
causal & $247/1000$ & $123/500$ & $1/754$ & $\sqrt{753/500}/(2\pi)$, $\Ctwoone$ & 24374032 \\
centered & n/a & n/a & n/a & $\Upre$ & 24372930 \\
\hline
\end{tabular}
\end{table}

\section{Sharpness: the exact scalar constant and the floor of the method}\label{sec:sharp}

A ceiling is only informative next to a lower bound. In the scalar case both are available,
and the interval between them is what the coupled constants of Sec.~\ref{sec:ceilings}
occupy.

\subsection{The scalar constant}

In the single-channel model with a normalized spectrum, the linear instability of the
isotropic linearized Eliashberg equation at coupling $\lambda$ is governed by a spectral top
$h$ of a self-adjoint positive operator built from the spectrum, and the onset is
$\tau=\Omega h$ in units of the phonon energy. For the one-parameter family of spectra with
a power-law weight at low frequency, indexed by $\gamma$, the spectral top is certified as
follows.

\begin{theorem}[scalar constant]\label{thm:scalar}
Write $c_\gamma$ for the scalar constant of the power-law family, so that the linear
condition reads $\kb\Tlin<c_\gamma\sqrt{g^{2}}$ with $g^{2}=\lambda\Omega^{2}$, and write
$\Cinf=c_{2}$. Then
\begin{equation}\label{eq:cinf}
0.1827262477723295078215179<\Cinf<0.1827262477723295176168720,
\end{equation}
an interval of relative width $5.4\times10^{-17}$ and of absolute width below
$9.8\times10^{-18}$, so $\Cinf$ rounds to $0.18272624777232952$ at seventeen decimal places.
The value is attained as the limit of the single Einstein line,
$\kb\Tlin/\sqrt{\lambda}\Omega\to\Cinf$ as $\lambda\to\infty$, so no constant smaller than
$\Cinf$ can hold for the scalar class, and the $\gamma=2$ normalisation is exactly the
strong-coupling Einstein ceiling. For the other members of the family reported by the
record, the enclosure of the same kind is available for every member, with relative width at
most $10^{-6}$ and endpoints rounded outward:
$c_{1/2}\in(0.80943835,0.80943861)$, $c_{1}\in(0.2558364,0.2558366)$,
$c_{3/2}\in(0.19999096,0.19999098)$ and $c_{3}\in(0.170998736,0.17099874)$. Only the
$\gamma=2$ member is certified to seventeen decimal places, and only that member is used as
a ceiling in this paper.
Status: \PROVED{} upstream, and \CERT{} for the enclosure, Ref.~\cite{gamma} with the
gauge-fixed eigenvalue normalisation of Ref.~\cite{finiteband}. A separate secondary
citation, unrefereed and not verified here, places the normalising limit in the interval
$1.318139<g(2)<1.31815$ \cite{kondratev}.
\end{theorem}

The proof has two halves. The upper half is an extremal principle for the spectrum at fixed
$\lambda$ and $M$: among all spectra with a given mass and second moment, the Einstein line
maximizes the spectral top, a statement proved by a transfer-by-transfer comparison and by
convexity in the spectrum in Ref.~\cite{rens} and used as an input in
Refs.~\cite{cutoff,gamma}. The lower half is an explicit evaluation of the Einstein limit,
in which the zero-transfer thermal term cancels between gap and self-energy and neither the
nonzero-transfer diagonal Hankel contribution nor the negative mass term is dropped
\cite{gamma}. The singular-value bounds that turn the comparison into an enclosure are the
content of Ref.~\cite{gamma}; what this paper imports is the resulting number and its
certified width, not the singular-value analysis, which is not reproduced here.

\begin{remark}[Prior art, stated plainly]\label{rem:prior}
The sharp scalar constant is not new here. It is the value of Conjecture~2 of
Kiessling, Altshuler and Yuzbashyan \cite{kay} as formulated in their programme, and an
unrefereed repository manuscript \cite{kondratev} states a proof of that bound. What
Refs.~\cite{gamma,finiteband} add is an independently certified enclosure, and what
Refs.~\cite{optimum,finiteenergy,anisosharp} add is the extension from one channel with
one spectrum to a coupled class of channels with different bounded densities of states and
a glue moment matrix. Any statement in this paper about a \emph{new} ceiling is a statement
about the coupled class.
\end{remark}

\subsection{The soft-glue limit: where the ceiling is attained}\label{sec:soft}

The scalar value $\Cinf$ belongs to the coupled problem, because the scalar case is a member
of the class and $\Cinf\le\sup\kb\Tlin/\sqrt\rho$. The reverse question, whether the supremum
is strictly larger, is open. What is proved is that the lower endpoint is attained, in a
limit in which the glue becomes soft and the channel structure is equitable; the statement is
a corollary of Ref.~\cite{finiteenergy}, and its ingredients are the angular-locking argument
and the exact limiting coefficient for scaled centered boxes recorded in the supplement of
Ref.~\cite{univglue}.

\begin{theorem}[soft-glue limit]\label{thm:softlimit}
Let $d\ge2$ channels carry counting measure, and let $\mathrm{d}\mathfrak m_{xy}$ be fixed
reciprocal entrywise nonnegative moment measures on positive energies, with
$J=\int\mathrm{d}\mathfrak m$ and $L=\int e^{-2}\mathrm{d}\mathfrak m$ finite and the graph of the positive
off-diagonal entries of $J$ connected. The same retarded interaction enters the normal and
the pairing equations, with no additional instantaneous term, and the branch is centered,
with even density of states and zero real shift. Channel $x$ carries the box density of
states $\mathbf 1_{|\xi|<W_x^{\varepsilon}}$ with $W_x^{\varepsilon}=\omega_x\varepsilon^{-p}$,
where $\omega_x>0$ and $p>1$ are fixed. Every boson energy $e$ is replaced by
$\sqrt\varepsilon\,e$ without changing its moment weight, so that $J$ is unchanged, and
$\rho=\|J\|$. Then
\begin{equation}\label{eq:softlimit}
\lim_{\varepsilon\downarrow0}\frac{\kb\Tlin^{\varepsilon}}{\sqrt\rho}
=\Cinf\sqrt{\frac{\mathbf 1^{\!\top}J\mathbf 1}{d\,\rho}},
\end{equation}
which equals $\Cinf$ if and only if $J\mathbf 1=\rho\mathbf 1$, and is strictly smaller
otherwise.
Status: \PROVED{}, Ref.~\cite{finiteenergy}, with the two ingredients taken from the
supplement of Ref.~\cite{univglue}.
\end{theorem}

Two explicit models show both outcomes. Two Lorentzian atoms at base energies $1$ and $2$
with
\begin{equation}\label{eq:atoms}
M_1=\begin{pmatrix}1/2&0\\0&0\end{pmatrix},\qquad
M_2=\begin{pmatrix}1/4&1/4\\1/4&3/4\end{pmatrix},
\end{equation}
give $J=M_1+M_2$ of norm one with equal row sums and $L=M_1+M_2/4$ with row sums $5/8$ and
$1/4$; with boxes $W=(1,2)/\varepsilon^{2}$ the leading normal widths satisfy
$\varepsilon s_{n,1}\to5/8$ and $\varepsilon s_{n,2}\to1/4$ at every fixed frequency, so the
spectra, the densities of states and the normal widths are all different, and yet
$\kb\Tlin\to\Cinf$. A single soft atom with $J=L=\bigl(\begin{smallmatrix}2&1\\1&1\end{smallmatrix}\bigr)$
has moment norm $(3+\sqrt5)/2$ and mean row sum $5/2$, so the limit is
$\Cinf\sqrt{5/(3+\sqrt5)}\approx0.977\,\Cinf$, strictly below $\Cinf$.

The limit is narrow in the following precise sense, and the disclaimer belongs to the
statement. It asserts neither equality nor excess of the finite-$\varepsilon$ ratio over
$\Cinf$; it is not a bandwidth bound uniform in $\varepsilon$; and it does not apply to
arbitrary centered bandwidths, to disconnected or degenerating moment graphs, or to
finite-energy optima. In particular it does not decide whether any single model at finite
glue attains $\Cinf$, and it does not decide whether the supremum over the class exceeds
$\Cinf$: it shows only that the interval $[\Cinf,\Cthree]$ cannot be shortened by moving its
lower endpoint upward.

\subsection{The floor of the frozen-sequence method}

The reserve argument has a limit, and the limit is proved.

\begin{proposition}[floor of the method]\label{prop:floor}
Let $\mathcal W_{\mathrm{cent}}$ be the set of weight sequences of the centered class. An
inequality of the form \eqref{eq:edge}, required to hold for every member of
$\mathcal W_{\mathrm{cent}}$ at every $t>0$ and every $r>0$ with a reserve $d$ independent
of $t$, $r$ and the weights, cannot have $d$ larger than $1-\Lambda/2$ where $\Lambda$ is
the largest value of $-I_{xy}/B_r$ over that set, since \eqref{eq:edge} reads
$-I_{xy}/B_r\le2(1-d)$. Explicit weight sequences and explicit trial vectors give
$\Lambda>1.357186$, hence $d<0.321407$ and no such proof gives a constant below
\begin{equation}\label{eq:floor}
\frac{\sqrt{\Lambda}}{2\pi}>0.185411 .
\end{equation}
Status: \PROVED{}, Ref.~\cite{anisosharp} (proved in Ref.~\cite{finiteenergy}).
\end{proposition}

This floor is a statement about the method and not a lower bound for the supremum over the
class: a proof that used the normal equation to tie the weights to the masses would not be
subject to it. The reserve $3189/10000=0.3189$ of Ref.~\cite{anisosharp}, the one that
yields the constant $\Cthree$, is below the reserve threshold $d^{*}=0.321407$ of the
floor argument, and the floor coefficient itself is $0.19\%$ below $\Cthree$. The floor sits at the fraction $0.887$ of the
interval $[\Cinf,\Cthree]$ measured from $\Cinf$, so at most the upper $11.3\%$ of that
interval is reachable by a sharper version of the same edge inequality. Whether the supremum equals
$\Cinf$ or lies strictly above it is \OPEN{}, and no conjecture is made here. A direct
finite-resolution sweep of $600$ admissible centered Poisson weight sequences at thirty-five
frequencies, on a grid in imbalance and energy, observed the smallest reserve $0.32147971$
at $r=1$, $t=10^{-4}$ and $b_0=c_0=0.92$; that observation is \NUM{}, is not an infimum, and
does not replace the analytic floor. The question is confined to several channels: for a
single channel, and for coupled channels with energy-independent densities of states, the
constant $\Cinf$ itself holds, so only channels with different energy-dependent profiles can
move the supremum above the scalar value.

\section{Coupling-resolved ceilings and ratios for real materials}\label{sec:material}

For a single isotropic channel the ceiling is a function of $\lambda$ at fixed $M$, written
$C(\lambda)$, because at fixed $\lambda$ and $M=\lambda\Omega^{2}$ the spectrum enters only
through the probability measure on $u=\omega^{2}/\Omega^{2}$ with $\int d\mu=1$ and
$\int u\,d\mu=1$. Strict growth of $C(\lambda)$ is a corollary of strict single-Einstein
optimality at fixed $(\lambda,M)$ \cite{rens}, and $C(\lambda)\uparrow\Cinf$ as
$\lambda\to\infty$.

The values are certified in exact arithmetic on a rational grid.

\begin{table}[h]
\caption{Coupling-resolved ceilings $C(\lambda)$ and the mean-square threshold for a
$300$ K ceiling. Columns two and three are the \CERT{} endpoints of the enclosure of
$C(\lambda)$, computed in exact rational arithmetic and checked by two independent programs;
the ratio $C/\Cinf$ and the threshold $M^{*}$ are \NUM{} evaluations of the printed
endpoints. Rows at $\lambda=2.5$ and $\lambda=6$ are omitted for space. The full set has
$19$ double-certified intervals: the nine grid couplings shown and omitted here, from
$\lambda=0.5$ to $10$, and the ten conventional and hydride couplings of
Sec.~\ref{sec:material}; all nineteen are tabulated in the supplement.\label{tab:clambda}}
\small
\begin{tabular}{@{}rrrrr@{}}
\hline
$\lambda$ & $C_{\mathrm{lo}}(\lambda)$ & $C_{\mathrm{hi}}(\lambda)$ & $C/\Cinf$ & $M^{*}$ at $300$ K (eV$^{2}$) \\
\hline
$0.5$ & $0.0557608245$ & $0.0557608803$ & $0.305161$ & $0.21494606$ \\
$1.0$ & $0.1146194547$ & $0.1146195694$ & $0.627274$ & $0.05087111$ \\
$1.5$ & $0.1381081893$ & $0.1381083275$ & $0.755821$ & $0.03503878$ \\
$2.0$ & $0.1497235932$ & $0.1497237430$ & $0.819388$ & $0.02981312$ \\
$3.0$ & $0.1610715064$ & $0.1610716675$ & $0.881491$ & $0.02576027$ \\
$4.0$ & $0.1666217589$ & $0.1666219257$ & $0.911866$ & $0.02407268$ \\
$10.0$ & $0.1763847604$ & $0.1763849369$ & $0.965296$ & $0.02148156$ \\
$\infty$ & $\Cinf$ & $\Cinf$ & $1$ & $0.02001641$ \\
\hline
\end{tabular}
\end{table}

The last column is the mean-square phonon energy, in electron volts squared, that a
$300$ K ceiling allows at that coupling. Since $\tau\le C(\lambda)\sqrt M$, a material whose
$\lambda$ and $M$ place it below that threshold, $M<M^{*}(300\,\mathrm{K})$, cannot have a
linear onset above $300$ K in this model class whatever its spectrum, and the threshold
decreases with $\lambda$ because $C(\lambda)$ increases.

\subsection{Ratios for four conventional and hydride cases}

The following table converts the ceiling into kelvin using literature couplings and
digitized second moments. It is \NUM{} in every entry that involves a measurement or a
digitization, and it is a ratio of a measured transition temperature to a computed
ceiling, not a validation of the theory and not a prediction for a material.

\begin{table}[h]
\caption{Measured transition temperatures against the computed ceilings. $M$ is
$\lambda\langle\omega^{2}\rangle$ evaluated from digitized Eliashberg functions. The
digitization uncertainty of $5$ to $10\%$ in $M$ stated in Sec.~\ref{sec:limits} is not a
confidence interval and is not propagated into the printed ratios, which also omit the
spread between published functions for the same material.
$T_{\mathrm{ceil}}(\lambda)$ uses the certified $C(\lambda)$ intervals of the supplement,
sampled in Table~\ref{tab:clambda},
$T_{\mathrm{ceil},\infty}$ uses $\Cinf$, and every ratio is \NUM{}, quoted to the digits
the underlying digitization supports; the certified intervals and the conversion script
are given in the supplement. For $\mathrm{MgB_{2}}$ the digitized reference spectrum is
isotropic while the material is strongly anisotropic, so its ratio is against an isotropic
reference only. For $\mathrm{LaH_{10}}$ the measured points and the density-functional
pressures do not coincide: the measured $244$ to $245$ K points are at $145$ and $135$ GPa,
paired here to a density-functional spectrum at $129$ GPa.\label{tab:mat}}
\small
\begin{tabular}{@{}lrrrrrr@{}}
\hline
material & $\lambda$ & $M$ (eV$^{2}$) & $T_{\mathrm{ceil}}(\lambda)$ (K) & $T_{\mathrm{ceil},\infty}$ (K) & $T_c$ (K) & $T_c/T_{\mathrm{ceil}}$ \\
\hline
Pb & $1.68$ & $0.00006249$ & $13.13$ & $16.76$ & $7.19$ & $0.548$ / $0.429$ \\
Nb & $1.26$ & $0.00034660$ & $27.90$ & $39.48$ & $9.25$ & $0.332$ / $0.234$ \\
$\mathrm{MgB_{2}}$ & $0.61$ & $0.0028654$ & $45.84$ & $113.51$ & $39$ & $0.851$ / $0.344$ \\
$\mathrm{LaH_{10}}$, $129$ GPa & $3.62$ & $0.036826$ & $367.18$ & $406.92$ & $244$--$245$ & $0.665$--$0.667$ / $0.600$--$0.602$ \\
\hline
\end{tabular}
\end{table}

The two ratios in the last column are $T_c/T_{\mathrm{ceil}}(\lambda)$ and
$T_c/T_{\mathrm{ceil},\infty}$. Reading the four rows: lead sits at $0.43$ of the scalar
ceiling and at $0.55$ of the ceiling that its own coupling allows; niobium at $0.23$ and
$0.33$; magnesium diboride at $0.34$ and $0.85$ in its isotropic reference; and
$\mathrm{LaH_{10}}$ at about $0.60$ of the scalar ceiling and $0.67$ of the ceiling at its
own coupling. Every measured value is below both ceilings, as it must be if the ceilings
are ceilings. What is informative is the size of the gap and its dependence on $\lambda$,
and the two columns order the four cases differently. Against the scalar ceiling the ratios
are $0.43$, $0.23$, $0.34$ and $0.60$; against the ceiling at the coupling of each material
they are $0.55$, $0.33$, $0.85$ and $0.66$ to $0.67$. The coupling-resolved ratio therefore
rises towards $\mathrm{MgB_2}$, whose coupling is the smallest in the table and whose
reference spectrum is isotropic, and it falls again for the hydride, whose coupling is the
largest. The scalar ceiling is not a tight statement at small $\lambda$, and the largest
ratio in the table belongs to the material whose reference spectrum is the least faithful,
so no ordering of materials by proximity to the ceiling is asserted. The only quantitative reading is the one the tables
support: every measured value lies strictly below the ceiling at its own coupling, and the
distance to the ceiling at fixed $300$ K is smallest where the mean square
$\lambda\langle\omega^{2}\rangle$ is largest.

\section{The ceiling is Coulomb robust; the optimizer is not}\label{sec:coulomb}

A ceiling in terms of $\rho$ alone would be vacuous if a pairing-only repulsion could be
used to raise the onset without limit, and a ceiling in terms of $C(\lambda)\sqrt M$ would be
vacuous if the maximizing spectrum changed whenever a repulsion was switched on. The
released record answers the two questions in opposite directions.

\begin{proposition}[Coulomb robustness of the ceiling]\label{prop:robust}
Let a pairing-only Coulomb repulsion act through a cutoff that is common to every spectrum
compared at fixed $\lambda$ and $M$, with pseudopotential $\muS\ge0$. The repulsion enters
the folded stability form as a positive semidefinite rank-one term, so no spectrum exceeds
$C(\lambda)\sqrt M$ at any $\muS\ge0$.
Status: \PROVED{}, Ref.~\cite{cutoff}.
\end{proposition}

\begin{proposition}[Fragility of the optimizer]\label{prop:fragile}
The optimizer is not robust. Three two-line spectra sharing $(\lambda,M,\muS,\omega_c)$ with
the Einstein line have a strictly larger linear onset, at $\muS/\lambda=2/3$, $3/4$ and $1$,
with temperature gains of at least $2.59$, $29.09$ and $82.57$ percent; and four further
counterexamples at $\muS/\lambda=0.4786$ to $0.4890$ and $\lambda=25.6$ to $28$ show that
$\muS/\lambda=1/2$ is not a threshold of the model. All are certified in exact rational
arithmetic by two independent verifiers.
Status: \CERT{}, Ref.~\cite{cutoff}.
\end{proposition}

The consequence for the present paper is a scoping statement. The ceiling of
Sec.~\ref{sec:material} survives any Coulomb repulsion; the identification of the Einstein
line as the worst case does not survive one, and the statement that a spectrum cannot beat
the ceiling is therefore weaker than the statement that a specific spectrum achieves it.

\section{Vertex corrections: a dichotomy}\label{sec:vertex}

Migdal's theorem makes the vertex a small correction when the phonon energy is small
compared with the electronic bandwidth. A ceiling that ignores vertices therefore has to
say what happens when they are retained. The released record treats the first crossed
skeleton, which is the leading vertex correction, in three regimes, and the answers differ
qualitatively. Throughout, $\chi=m\sqrt\lambda$ is the crossing parameter, with
$m=\Omega_{\max}/W$ the mass ratio, $\theta=\Omega_{\max}/(2\nu)=(2\pi\eta\Cinf\sqrt\lambda)^{-1}$
the Matsubara ratio at $\tau=\eta\Tceil$, $x=\pi\tau/W=\pi\eta\Cinf\sqrt\lambda\,m$ the
adiabatic parameter, $A$ the density-of-states ratio, and $B=\max(1,A v_{\max}R)$ a
geometric factor. The crossing parameter is written $\chi$ in this section so that it is
never confused with the moment norm $\rho$ of Sec.~\ref{sec:ceilings}.

\begin{theorem}[flat density of states: an inflated ceiling]\label{thm:flatceiling}
Take a flat density of states on $[-W,W]$, the local first crossed skeleton and even-frequency
singlet pairing. For any centered particle-hole-symmetric normal solution with positive
reflected widths,
\begin{equation}\label{eq:V1}
\kb\Tlin\le\Cinf\sqrt{AM}\Bigl[1+30\bigl((1+\lambda)\chi+\chi^{2}\bigr)\Bigr].
\end{equation}
The bound requires neither a regular branch nor an existence guard. If
$(1+\lambda)\chi\le1/100$, a centered positive-width solution exists pointwise for every
$\tau\in[\Cinf\sqrt M,2\Cinf\sqrt M]$, and at fixed finite $\lambda$ the uninflated ceiling
holds whenever $\epsilon=30[(1+\lambda)\chi+\chi^{2}]<\delta_\lambda$, with
$\delta_\lambda=1-\lambda h((2\pi\Cinf\sqrt\lambda)^{-1})>0$ and $h$ the scalar Einstein
spectral top per unit coupling.
Status: \PROVED{}, Ref.~\cite{msvertex}.
\end{theorem}

The last clause is a statement about the direct linear threshold of the same branch, not a
proof that the vertex-free inequality $\kb\Tlin\le\Cinf\sqrt\rho$ survives the vertex. What
the record proves is the inflated inequality \eqref{eq:V1} for every centered
particle-hole-symmetric normal solution with positive reflected widths, together with the
$\epsilon<\delta_\lambda$ regime in which the inflation is absorbed into the exact special
flow. The inflation is first order in the crossing parameter $\chi$ and linear in
$\lambda$, so the
vertex-free ceiling of Sec.~\ref{sec:ceilings} holds by continuity in the window
$\chi\ll1$; it is not uniform in $\lambda$, and no uniformity in $\lambda$ is claimed.

\begin{theorem}[nested sheets and curved surfaces]\label{thm:dichotomy}
Fix $\lambda>0$ and $\eta>0$, and set $\tau=\eta\Tceil$, $\nu=\pi\tau$,
$m=\Omega_{\max}/W$ and $y=\xi/\nu$ on each sheet, so that
$\tau/\Omega_{\max}=\eta\Cinf\sqrt\lambda$ and $\Omega_{\max}/(2\nu)=\theta$ depend on $m$
only through $\lambda$. On two perfectly nested one-dimensional sheets, that is
$k=\pm(k_F+\frac{\pi\nu}{2W}y)$, with $4k_F\notin2\pi\mathbb{Z}$, the crossed self-energy has
a finite limit as $m\to0$ at every fixed $y\ne0$ and every fixed $\eta$: of the eight sign
channels three keep the middle line in the collar, their prefactors are free of $m$, the
real rainbow shift is a translation, and the remaining two integrals are exact Cauchy
residues whose value at fixed $y$ is independent of $m$. No estimate of the form $O(m)$ for
the crossed self-energy can hold at fixed $\eta$. On Fermi surfaces in $D\ge2$ dimensions
whose level sets in the collar are strictly convex with principal curvatures bounded below,
the crossed pairing correction obeys, for a sign-preserving normal branch with
$u_n(\mathbf k)\ge c_0|\omega_n|$, $c_0\in(0,1]$, and bounded
$|\operatorname{Re}\Sigma_n(\mathbf k)|\le\Delta_R$,
\begin{equation}\label{eq:coarea}
|F_X[y]|\le10L_XQ_0[y],\qquad
L_X\le\frac{\gamma^{3}c_0^{\alpha-1}\CLC}{\pi^{2}}B^{2}x^{\alpha}
\Bigl(1+\log_+\frac{1}{c_0x}\Bigr)^{3+\beta},
\end{equation}
uniformly in the shape of the level sets, with $\gamma=1+\Delta_R/(c_0\nu)$,
$\alpha=1$, $\beta=1$ in $D=2$ and $\alpha=1$, $\beta=0$ in $D\ge3$. Here $\CLC$ is the
level-crossing constant of the record, assembled from the geometric data of the collar, and
it enters only through the prefactor. In the small-shift regime $\Delta_R\le c_1\nu$ one has
$\gamma=O(1)$; in the wider regime $\Delta_R\le C_{\mathrm{Re}}\nu(1+\log_+(1/x))$ the factor
$\gamma^{3}$ adds at most three powers of $1+\log_+(1/(c_0x))$. The first crossed skeleton
therefore shifts the linear onset of that branch by a relative factor at most
\begin{equation}\label{eq:shift}
1+O\bigl((A v_{\max}R)^{2}x^{\alpha}(1+\log_+(1/x))^{3+\beta}\bigr),\qquad x\to0,
\end{equation}
in the small-shift regime $\Delta_R\le c_1\nu$, and at most
\begin{equation}\label{eq:shiftwide}
1+O\bigl((A v_{\max}R)^{2}x^{\alpha}(1+\log_+(1/x))^{6+\beta}\bigr),\qquad x\to0,
\end{equation}
in the wider regime $\Delta_R\le C_{\mathrm{Re}}\nu(1+\log_+(1/x))$, at each fixed $\eta$ and
each fixed equation-of-state family. That is linear in the adiabatic ratio up to a
logarithm cubed in \eqref{eq:shift} and up to a logarithm to the sixth power in
\eqref{eq:shiftwide}, with the corresponding powers one higher in $D=2$, where $\beta=1$:
a logarithm to the fourth and to the seventh. Since $x=\pi\eta\Cinf\sqrt\lambda\,m$ and
$\chi=m\sqrt\lambda$ fix each other at fixed $\eta$ and $\lambda$, both shifts are the same
statement in $\chi$.
Status for this part: the vanishing is unconditional in the shape of the level sets within
the stated curvature class. Separately, in an inversion-invariant projected continuum domain with constant
spatial transfer and even marginal density of states, a Young-ball contraction gives an
all-$x$ regime without any curvature hypothesis: for $A\lambda<1/9$,
\begin{equation}\label{eq:young}
\kb\Tlin\le\tau_Y=\frac{9AM_1}{2(1-9A\lambda)},\qquad M_1=\int\Omega\,d\lambda(\Omega),
\end{equation}
and for $A\lambda\le1/26^{2}$ one has $\tau_Y\le(117/667)\sqrt{AM}<\Cinf\sqrt{AM}$ with no
adiabatic assumption.
Status: \PROVED{} with the scope stated in each part, Ref.~\cite{msvertex}. Both extra
hypotheses of the all-$x$ part are needed: without a particle-hole-even density of states
and a constant transfer, the bare rainbow shift
$\mathrm{Re}\,\Sigma^{(0)}_{\mathrm{RB}}=\Theta(\nu A v_{\max}R\log(1/x))$ leaves every ball
of fixed radius as $x\to0$, however small $AR$ is.
\end{theorem}

The dichotomy is the physical content. Nested sheets, such as the $\gamma$-pocket sheets of
a layered material near a van Hove filling, make the crossed term survive the adiabatic
limit, so no universal vertex-free ceiling can exist for them; curved surfaces make it
vanish linearly in the adiabatic ratio up to a logarithm cubed, a logarithm to the fourth
power in two dimensions, so the ceiling
of Sec.~\ref{sec:ceilings} is stable there in the stated window. The all-coupling,
all-momentum, fully crossed statement is \OPEN{}, and it is the largest open problem of the
programme: the universal onset statement of the first version of Ref.~\cite{msvertex} is not
settled by Theorem~\ref{thm:dichotomy}. What is \PROVED{} beyond the ME ceiling is that a
termwise absolute triangle estimate cannot work (\cite{msvertex}, the two failing estimates:
at fixed positive $\lambda$ and $\tau=\Omega$ the dynamic triangle at zero spatial transfer
tends to $\lambda$ for bare and to $\lambda/(1+\lambda)$ for dressed flat-band lines, so an
$O(\lambda m)$ essential-supremum estimate is unavailable).

\section{Comparison with the literature}\label{sec:lit}

\subsection{Allen and Dynes}

The Allen--Dynes analysis \cite{allendynes} is a fit to analytic and numerical solutions of
the Eliashberg equations, and its abstract quotes a large-coupling coefficient $0.15$ in
units of $\sqrt{\lambda\langle\omega^{2}\rangle}$ at $\muS=0.1$. The exact scalar
coefficient at $\muS=0$ is $\Cinf=0.1827262478$, which is $21.8\%$ above that coefficient.
The comparison is therefore between a bound at zero repulsion and a fit at
$\muS=0.1$. Since
$\muS>0$ can only reduce the onset in the model class where the ceiling is Coulomb robust
(Proposition~\ref{prop:robust}), the comparison is not an inconsistency: it says that the
Allen--Dynes strong-coupling coefficient is strictly conservative with respect to a sharp
ceiling at zero repulsion, and that the empirical rule and the bound agree in magnitude and
sign. We do not convert this into material-specific predictions, because the fit
coefficients of the Allen--Dynes formula could not be verified from a primary source for
this manuscript.

\subsection{Moussa and Cohen}

Moussa and Cohen \cite{moussa} state two bounds. Their Eq.~(10) writes
$\kb T_c\le f(\lambda,\muS)\sqrt{\langle\omega^{2}\rangle}$ with $f$ given by an empirical
fit whose maximum deviation is $0.03$, so $f$ is a fit and not a bound. Their Eq.~(7) gives
the fixed-moment limit $0\le\kb T_c\le0.183\sqrt M$. The constant proved and certified here
is $0.1827262478$, which is $0.15\%$ below $0.183$: the present result replaces their
decimal by a certified transcendental value and extends the statement from one channel
with one spectrum to a coupled class with channel-dependent bounded densities of states.
Our reading of their Eqs.~(7), (10) and (11) relies on a secondary summary, which records
them; their full text was not re-read for this manuscript, and we attach the label
\KNOWN{} rather than \PROVED{} to the comparison.

\subsection{The Kiessling--Altshuler--Yuzbashyan programme}

The programme formulated in \cite{kay} asks for bounds on $T_c$ in Eliashberg theory in terms
of spectral moments. The scalar constant $\Cinf$ is the value that closes Conjecture~2 of
that programme, and an unrefereed repository manuscript \cite{kondratev} states a proof of
it, with an exact-arithmetic eigenvalue enclosure $1.318139<g(2)<1.31815$ that contains our
certified value. The contribution of the coupled records cited here is therefore
twofold: an independent certified enclosure of that number, and the extension to the
multichannel problem, where the constant is not sharp: the certified supremum bound is
$\Cthree$ with $\Cthree/\Cinf-1=1.658\%$, and the sharpest strict constant of the same record
is $\Cstar$, $1.725\%$ above $\Cinf$ and above $\Cthree$.

\subsection{Vertex and nonadiabatic literature}

Vertex corrections to Eliashberg theory have been studied as a practical correction for a
long time, both locally conserving \cite{miller} and by local vertex computations at
half filling \cite{paci}. The results above are of a different character: they are bounds on
how much the first crossed skeleton can move the linear onset, in a geometry-dependent
dichotomy, and they do not address the self-consistent crossed operator.

\section{Limits: what is not claimed}\label{sec:limits}

The bounds are narrow statements and the following list is part of their content.

\begin{enumerate}
\item[\rm 1.] \emph{Linear, not nonlinear.} The bound is on the linear onset of the folded
stability form of linearized Eliashberg theory, tested against a fixed positive normal
branch. It says nothing about the nonlinear gap equation at low temperature, and $\Tlin$ is
not the transition temperature of any exact Hamiltonian. No statement here identifies
$\Tlin$ with a measured $T_c$.
\item[\rm 2.] \emph{A model class, not a material.} The glue is retarded and the same
nonnegative transfer enters the normal and the pairing equations. Magnetic order, competing
orders, anisotropic order parameters beyond the channel structure of
Sec.~\ref{sec:model}, and off-diagonal $\eta$-pairing are outside. $\Cinf$ and the coupled
constants are properties of this class and are not bounds for any particular compound.
\item[\rm 3.] \emph{Vertices only in the regimes stated.} The fully crossed self-consistent
vertex operator is not covered. Theorem~\ref{thm:flatceiling} inflates the ceiling by a
term first order in the crossing parameter $\chi$ and linear in $\lambda$, and shows that
the uninflated inequality there concerns the direct linear threshold of the same branch,
Theorem~\ref{thm:dichotomy} shows that the crossed correction cannot vanish with the phonon
mass ratio on nested one-dimensional sheets and vanishes linearly in the adiabatic ratio up
to a logarithm cubed, a logarithm to the fourth power in two dimensions, on the curved
surfaces of the stated curvature class, and the Young-ball regime gives an
all-$x$ statement only for
$A\lambda<1/9$. Outside these regimes the ceiling of Sec.~\ref{sec:ceilings} is a
mean-field-with-rainbow-vertex statement.
\item[\rm 4.] \emph{The constant is not sharp in the coupled class.} The chain shows
$\Cinf\le\sup\kb\Tlin/\sqrt\rho\le\Cthree$ with $\Cthree/\Cinf-1=1.658\%$, and
Proposition~\ref{prop:floor} shows that no frozen-sequence edge argument can close more than
the upper $11.3\%$ of that interval. Which end is correct is \OPEN{}. The optimal coupled
coefficient and the full shifted-versus-centered comparison are also \OPEN{}.
\item[\rm 5.] \emph{The scalar sharpness is not claimed as new.} Theorem~\ref{thm:scalar}
is imported; the prior proof is an unrefereed repository manuscript, and we state that
plainly in Remark~\ref{rem:prior}.
\item[\rm 6.] \emph{Nothing is measured.} The four material rows of
Table~\ref{tab:mat} are ratios whose denominators are computed and whose numerators are
literature values. The second moments are digitized, with $5$ to $10$ percent uncertainty,
the pressure of a density-functional spectrum need not equal the pressure of a measurement,
and $\mathrm{MgB_{2}}$ is compared against an isotropic reference while it is strongly
anisotropic. No material is claimed to be a superconductor by these bounds, and no claim of
patentability, priority or novelty beyond the records cited is made.
\item[\rm 7.] \emph{Formalization is partial.} The finite rational core of the reserve
inequalities was checked in Lean~4.35.0 in the released records. The analytic part (Harnack
rearrangement, Poisson comparison, the width cone, the reduction of
Sec.~\ref{sec:reserve}) is human-reviewed mathematics and is not formalized.
\item[\rm 8.] \emph{The lattice layer is out of scope.} Pair order in the slow-phonon
Holstein model, its suppression for slow lattice phonons \cite{slowphonon}, and the
curvature of defect-insertion bounds \cite{limits} are a separate programme, mapped in
Ref.~\cite{capstone}. Nothing in the present paper is a bound on a lattice pair temperature.
\end{enumerate}

\section{Reproducibility, data availability and disclosure}\label{sec:repro}

Every finite statement of Sec.~\ref{sec:finite} is reproduced by exact-arithmetic programs
that are archived with the corresponding release, together with the guarded mutant runs and,
where present, the Lean sources: the replay supplements of Refs.~\cite{univglue,optimum,
finiteenergy,anisosharp,msvertex}. The scalar enclosure of Theorem~\ref{thm:scalar} is
reproduced by the replay of Ref.~\cite{gamma}, which exits successfully including the invalid
upper-bound injection and the seventeen-place rounding check. The coupling-resolved ceiling
$C(\lambda)$ of Table~\ref{tab:clambda} is computed for this paper: its exact rational
endpoints were produced by two independently written programs working in exact arithmetic,
and both programs, their mutation tests and the outward-rounded endpoints ship as
supplementary material with this submission. No data set is new here. The couplings and
second moments of Table~\ref{tab:mat} are evaluated in this paper from digitized Eliashberg
functions of the literature, and the transition temperatures are literature values; the
digitization records, with the figure and table each value was read from, and the script that
converts them into Table~\ref{tab:mat} also ship as supplementary material. The digitized
Eliashberg functions themselves are figures in the cited papers and are not redistributed.
The source of this manuscript, its build record and these files are deposited together:
the preprint is DOI \DOI{\PAPERDOI}, and the supplementary archive, with the repository
that carries it, is DOI \DOI{\SUPPLEMENTDOI}, the source and replay scripts being at
\url{https://github.com/SuperMatics-Inc/migdal-eliashberg-pairing-onset-ceilings}.

\emph{Use of AI assistance.} The author used an AI assistant for code generation, for exact
rational and interval arithmetic checks, for LaTeX production and for editorial work on the
text. All mathematical statements, hypotheses and constants were reviewed by the author, and
every numerical constant in this paper is traceable either to a released record whose
certification is described in Sec.~\ref{sec:finite} or to the supplementary verification
package of Sec.~\ref{sec:repro}. No claim in the paper rests on an unchecked model output.

\begin{acknowledgments}
The author thanks the maintainers of the exact-arithmetic and proof-assistant toolchains
used by the cited records. No external funding was received for this work.
\end{acknowledgments}

\bibliography{refs}

\end{document}
